% ================================================================== Chapitre 3
\chapter{Grandeurs de réaction et grandeurs standard}
\label{ch:grandeurs}

On considère un système siège d'\emph{une seule} réaction $0 = \sum_i\nu_i\mathrm{A}_i$, d'avancement $\xi$. Toute grandeur extensive $X$ du système ($G$, $H$, $S$, $V$\dots) dépend alors de $T$, $P$ et $\xi$. La question centrale de ce chapitre : \emph{comment $X$ varie-t-elle quand la réaction avance ?}

\section{Grandeurs de réaction}

\begin{definition}[grandeur molaire partielle]
La \textbf{grandeur molaire partielle} de $\mathrm{A}_i$ associée à $X$ est $X_{m,i} = \left(\dfrac{\partial X}{\partial n_i}\right)_{T,P,n_{j\neq i}}$. Pour $X = G$, c'est le potentiel chimique : $G_{m,i} = \mu_i$.
\end{definition}

\begin{definition}[grandeur de réaction]
La \textbf{grandeur de réaction} associée à $X$ est la dérivée de $X$ par rapport à l'avancement, à $T$ et $P$ fixées :
\[ \boxed{\Delta_r X = \left(\frac{\partial X}{\partial \xi}\right)_{T,P} = \sum_i \nu_i\,X_{m,i}} \qquad\text{en particulier}\qquad \boxed{\Delta_r G = \sum_i \nu_i\,\mu_i} \]
Elle s'exprime en $\mathrm{J\cdot mol^{-1}}$ (ou $\mathrm{J\cdot K^{-1}\cdot mol^{-1}}$ pour $\Delta_rS$) : « par mole d'avancement ». Elle dépend de $T$, de $P$ \emph{et de la composition} du système.
\end{definition}

\begin{demo}
À $T$ et $P$ fixées, $\dd X = \sum_i X_{m,i}\,\dd n_i = \sum_i X_{m,i}\,\nu_i\,\dd\xi$, d'où $\left(\partial X/\partial\xi\right)_{T,P} = \sum_i\nu_iX_{m,i}$.
\end{demo}

Le symbole $\Delta_r$ n'est \emph{pas} une différence entre deux états : c'est un opérateur de dérivation (« opérateur de Lewis »). Une grandeur de réaction est une \emph{pente} : $\Delta_rG$ est la pente de la courbe $G(\xi)$ (chapitre~\ref{ch:evolution}).

\begin{pieges}
\begin{itemize}
  \item Une grandeur de réaction est attachée à une \emph{équation écrite} : si l'on multiplie tous les coefficients par 2, $\Delta_rX$ est multipliée par 2. Toujours préciser l'équation utilisée.
  \item Ne pas confondre $\Delta_rH$ (en $\mathrm{J\cdot mol^{-1}}$, intensive) et $\Delta H$ (en J, variation entre deux états). Le lien est $\Delta H = \int \Delta_rH\,\dd\xi$.
\end{itemize}
\end{pieges}

\section{Grandeurs standard de réaction}

\begin{definition}[grandeur standard de réaction]
La \textbf{grandeur standard de réaction} est la grandeur de réaction calculée en supposant \emph{chaque constituant dans son état standard} à la température $T$ :
\[ \boxed{\Delta_r X^\circ(T) = \sum_i \nu_i\,X^\circ_{m,i}(T)} \]
en particulier
\[ \Delta_rG^\circ = \sum_i\nu_i\,\mu_i^\circ,\qquad \Delta_rH^\circ = \sum_i\nu_iH^\circ_{m,i},\qquad \Delta_rS^\circ = \sum_i\nu_iS^\circ_{m,i}. \]
Les $X^\circ_{m,i}$ sont les grandeurs molaires des constituants purs dans leur état standard : $\Delta_rX^\circ$ \emph{ne dépend que de $T$}.
\end{definition}

On peut se représenter $\Delta_rX^\circ$ comme la variation de $X$ pour une transformation \emph{fictive} : une mole d'avancement, à $T$ et $P^\circ$ fixées, où réactifs et produits resteraient chacun dans leur état standard (donc purs et séparés). Une telle réaction « totale fictive » n'est pas réalisable, mais elle définit une référence tabulable.

\begin{figure}[H]
\centering
\begin{tikzpicture}[>=Stealth,box/.style={draw=navy,line width=0.8pt,rounded corners=2pt,fill=cdefbg,minimum width=3.3cm,minimum height=1cm,align=center,font=\small}]
  \node[box] (R) at (0,0) {Réactifs\\dans l'état standard};
  \node[box] (P) at (7.4,0) {Produits\\dans l'état standard};
  \draw[->,line width=1.1pt,c3] (R) -- node[above,font=\small,black] {réaction fictive, totale} node[below,font=\small,black] {$T$ fixée, $P = P^\circ$} (P);
  \node[font=\small,c3] at (3.7,-1.05) {$\Delta X = \Delta_rX^\circ$ pour $\xi = \SI{1}{mol}$};
\end{tikzpicture}
\caption{Interprétation d'une grandeur standard de réaction.}
\end{figure}

\begin{resultat}{Relation fondamentale}
\[ \boxed{\Delta_rG^\circ(T) = \Delta_rH^\circ(T) - T\,\Delta_rS^\circ(T)} \]
\end{resultat}

\begin{demo}
Pour chaque constituant dans son état standard, $\mu_i^\circ = G^\circ_{m,i} = H^\circ_{m,i} - TS^\circ_{m,i}$. On multiplie par $\nu_i$ et on somme.
\end{demo}

\section{Signification physique et signe}

\begin{definition}[réaction endothermique, exothermique, athermique]
Une réaction est dite \textbf{exothermique} si $\Delta_rH^\circ < 0$, \textbf{endothermique} si $\Delta_rH^\circ > 0$, \textbf{athermique} si $\Delta_rH^\circ \approx 0$. Le sens physique (le système cède ou reçoit de la chaleur) est établi au chapitre~\ref{ch:premier}.
\end{definition}

\begin{resultat}{Signe de l'entropie standard de réaction}
L'entropie molaire d'un gaz (de l'ordre de $\SI{200}{J.K^{-1}.mol^{-1}}$) est très supérieure à celle d'un solide ou d'un liquide (quelques dizaines). En notant $\Delta\nu_{\text{gaz}} = \sum_{i\,\text{gaz}}\nu_i$ :
\[ \Delta\nu_{\text{gaz}} > 0 \;\Rightarrow\; \Delta_rS^\circ > 0, \qquad \Delta\nu_{\text{gaz}} < 0 \;\Rightarrow\; \Delta_rS^\circ < 0, \qquad \Delta\nu_{\text{gaz}} = 0 \;\Rightarrow\; |\Delta_rS^\circ| \text{ faible}. \]
Interprétation : l'entropie mesure le désordre ; créer du gaz augmente fortement le nombre de configurations accessibles.
\end{resultat}

\begin{exemple}[synthèse de l'ammoniac]
\ce{N2(g) + 3 H2(g) = 2 NH3(g)} : $\Delta\nu_{\text{gaz}} = 2 - 4 = -2 < 0$, donc on prévoit $\Delta_rS^\circ < 0$. Avec les entropies molaires standard (tableau~\ref{tab:donnees}) :
\begin{align*} \Delta_rS^\circ &= 2S^\circ_m(\ce{NH3}) - S^\circ_m(\ce{N2}) - 3S^\circ_m(\ce{H2})\\ &= 2(192{,}8) - 191{,}6 - 3(130{,}7) = \SI{-198,1}{J.K^{-1}.mol^{-1}}. \end{align*}
\end{exemple}

\section{Grandeurs standard de formation et loi de Hess}

\begin{definition}[réaction de formation, grandeur standard de formation]
La \textbf{réaction standard de formation} d'un constituant à la température $T$ est la réaction qui forme \emph{une mole} de ce constituant, dans son état standard, à partir des éléments qui le composent pris dans leur \emph{état standard de référence}.

\noindent La \textbf{grandeur standard de formation} $\Delta_fX^\circ$ est la grandeur standard de cette réaction.
\end{definition}

\begin{exemple}
Formation de l'ammoniac : $\tfrac12\,\ce{N2(g)} + \tfrac32\,\ce{H2(g)} \longrightarrow \ce{NH3(g)}$, avec $\Delta_fH^\circ(\ce{NH3(g)}) = \SI{-45,9}{kJ.mol^{-1}}$ (donnée de table).

\noindent Formation de l'ozone : $\tfrac32\,\ce{O2(g)} \longrightarrow \ce{O3(g)}$, donc $\Delta_fX^\circ(\ce{O3}) \neq 0$ alors que c'est un corps simple.
\end{exemple}

\begin{resultat}{Conséquence immédiate}
\[ \boxed{\Delta_fX^\circ(\text{corps simple dans son ESR}) = 0} \]
car sa « réaction de formation » ne transforme rien (\ce{O2(g) -> O2(g)}). En revanche $\Delta_fX^\circ \neq 0$ pour un corps composé \emph{ou} pour un corps simple hors de son ESR (\ce{O3}, diamant, \ce{Br2(g)} à $\SI{298}{K}$\dots).
\end{resultat}

\begin{resultat}{Loi de Hess}
\[ \boxed{\Delta_rH^\circ = \sum_i \nu_i\,\Delta_fH^\circ_i} \qquad\text{et de même}\qquad \Delta_rG^\circ = \sum_i\nu_i\,\Delta_fG^\circ_i \]
\end{resultat}

\begin{demo}
$H$ est une fonction d'état : sa variation ne dépend pas du chemin suivi. On passe des réactifs aux produits en deux étapes fictives : (1) décomposer les réactifs en éléments dans leur ESR, ce qui contribue $-\sum_{\text{réactifs}}|\nu_i|\,\Delta_fH^\circ_i$ ; (2) reformer les produits à partir de ces éléments, ce qui contribue $+\sum_{\text{produits}}\nu_i\,\Delta_fH^\circ_i$. Avec les coefficients algébriques, les deux sommes n'en font qu'une.
\end{demo}

\begin{figure}[H]
\centering
\begin{tikzpicture}[>=Stealth,box/.style={draw=navy,line width=0.8pt,rounded corners=2pt,fill=cdefbg,inner sep=5pt,align=center,font=\small}]
  \node[box] (R) at (0,2.6) {\ce{CH4(g) + Cl2(g)}};
  \node[box] (P) at (8.6,2.6) {\ce{CH3Cl(g) + HCl(g)}};
  \node[box,fill=cmetbg,draw=c4] (E) at (4.3,0) {\ce{C(graphite) + 2 H2(g) + Cl2(g)}\\{\scriptsize éléments dans leur ESR}};
  \draw[->,line width=1.2pt,c3] (R) -- node[above,font=\small] {$\Delta_rH^\circ = \;?$} (P);
  \draw[->,line width=1pt,navy] (R.south) -- node[below left,font=\small,align=right,pos=0.5] {$-\Delta_fH^\circ(\ce{CH4})$\\$-\,\Delta_fH^\circ(\ce{Cl2}) = 0$} (E.north west);
  \draw[->,line width=1pt,navy] (E.north east) -- node[below right,font=\small,align=left,pos=0.5] {$\Delta_fH^\circ(\ce{CH3Cl})$\\$+\,\Delta_fH^\circ(\ce{HCl})$} (P.south);
\end{tikzpicture}
\caption{Cycle de Hess pour la chloration du méthane : on passe par les éléments dans leur état standard de référence.}
\label{fig:hess}
\end{figure}

\begin{exemple}[chloration du méthane]
\ce{CH4(g) + Cl2(g) -> CH3Cl(g) + HCl(g)}. D'après la loi de Hess et les données à $\SI{298}{K}$ :
\begin{align*} \Delta_rH^\circ &= \Delta_fH^\circ(\ce{CH3Cl}) + \Delta_fH^\circ(\ce{HCl}) - \Delta_fH^\circ(\ce{CH4}) - \underbrace{\Delta_fH^\circ(\ce{Cl2})}_{0}\\ &= -81{,}9 - 92{,}3 + 74{,}6 = \SI{-99,6}{kJ.mol^{-1}}. \end{align*}
La réaction est exothermique.
\end{exemple}

\begin{resultat}{Et pour l'entropie ?}
Les tables ne donnent pas $\Delta_fS^\circ$ mais les \textbf{entropies molaires standard absolues} $S^\circ_m$ (le troisième principe fixe $S_m = 0$ pour un cristal parfait à $\SI{0}{K}$). On calcule donc directement
\[ \Delta_rS^\circ = \sum_i\nu_i\,S^\circ_{m,i}. \]
\end{resultat}

\begin{pieges}
\begin{itemize}
  \item $\Delta_fH^\circ(\ce{O2}) = 0$ \emph{mais} $S^\circ_m(\ce{O2}) = \SI{205,2}{J.K^{-1}.mol^{-1}} \neq 0$ : les entropies des corps simples ne sont \emph{jamais} nulles. Les oublier dans $\Delta_rS^\circ$ est l'erreur la plus fréquente.
  \item Unités : $\Delta_rH^\circ$ en $\mathrm{kJ\cdot mol^{-1}}$, $\Delta_rS^\circ$ en $\mathrm{J\cdot K^{-1}\cdot mol^{-1}}$. Dans $\Delta_rG^\circ = \Delta_rH^\circ - T\Delta_rS^\circ$, convertir avant de soustraire.
  \item Bien vérifier l'état physique : $\Delta_fH^\circ(\ce{H2O(l)}) = -285{,}8$ et $\Delta_fH^\circ(\ce{H2O(g)}) = \SI{-241,8}{kJ.mol^{-1}}$ ; la différence est l'enthalpie de vaporisation.
\end{itemize}
\end{pieges}

\section{Combinaisons linéaires de réactions}

\begin{resultat}{Combinaison linéaire}
Si l'équation (3) s'obtient comme $(3) = a\,(1) + b\,(2)$, alors
\[ \Delta_rX_3^\circ = a\,\Delta_rX_1^\circ + b\,\Delta_rX_2^\circ \quad (X = H, S, G) \qquad\text{et}\qquad K_3^\circ = \left(K_1^\circ\right)^a\left(K_2^\circ\right)^b. \]
\end{resultat}

\begin{demo}
$\Delta_rX^\circ = \sum_i\nu_iX^\circ_{m,i}$ est linéaire en les coefficients $\nu_i$, d'où la première relation. Pour $K^\circ$, on utilise $\Delta_rG^\circ = -RT\ln K^\circ$ (établi au chapitre~\ref{ch:evolution}) : $\ln K_3^\circ = a\ln K_1^\circ + b\ln K_2^\circ$.
\end{demo}

\begin{exemple}[enthalpie de formation de CO]
On ne peut pas mesurer directement \ce{C + 1/2 O2 -> CO} (il se forme toujours un peu de \ce{CO2}). On combine :
\begin{align*} (1)&\quad \ce{C(gr) + O2 -> CO2} & \Delta_rH_1^\circ &= \SI{-393,5}{kJ.mol^{-1}}\\ (2)&\quad \ce{CO + 1/2 O2 -> CO2} & \Delta_rH_2^\circ &= \SI{-283,0}{kJ.mol^{-1}} \end{align*}
$(3) = (1) - (2)$ donne \ce{C(gr) + 1/2 O2 -> CO}, d'où $\Delta_fH^\circ(\ce{CO}) = -393{,}5 + 283{,}0 = \SI{-110,5}{kJ.mol^{-1}}$.
\end{exemple}

\section{Influence de la température : approximation d'Ellingham}

\begin{definition}[approximation d'Ellingham]
On suppose que $\Delta_rH^\circ$ et $\Delta_rS^\circ$ sont \textbf{indépendantes de la température} sur tout intervalle où aucun constituant ne change d'état. Alors
\[ \boxed{\Delta_rG^\circ(T) = \Delta_rH^\circ - T\,\Delta_rS^\circ} \]
est une \textbf{fonction affine} de $T$, de pente $-\Delta_rS^\circ$ et d'ordonnée à l'origine $\Delta_rH^\circ$.
\end{definition}

Cette approximation est imposée par le programme. Elle est justifiée car $\Delta_rH^\circ$ et $\Delta_rS^\circ$ varient peu avec $T$ (leurs dérivées font intervenir $\Delta_rC_p^\circ$, petite devant $\Delta_rH^\circ/T$), alors que le terme $T\Delta_rS^\circ$ change beaucoup.

\begin{figure}[H]
\centering
\begin{tikzpicture}
\begin{axis}[cours,xlabel={$T$ (K)},ylabel={$\Delta_rG^\circ$ ($\mathrm{kJ\cdot mol^{-1}}$)},
  xmin=300,xmax=900,ymin=-40,ymax=100,
  xtick={300,400,500,600,700,800,900}]
  \fill[c4!12] (axis cs:300,0) -- (axis cs:300,-32.37) -- (axis cs:463.4,0) -- cycle;
  \fill[c3!10] (axis cs:463.4,0) -- (axis cs:900,86.49) -- (axis cs:900,0) -- cycle;
  \draw[black,line width=0.6pt] (axis cs:300,0) -- (axis cs:900,0);
  \addplot[c1] table[x=T,y=DrG_kJ]{data/nh3_drg.dat};
  \addplot[c3,only marks,mark=*,mark size=2.2pt] coordinates {(463,0)};
  \node[c3,font=\small,anchor=south east] at (axis cs:458,3) {$T_i = \SI{463}{K}$};
  \node[c4,font=\scriptsize] at (axis cs:340,-8) {$K^\circ>1$};
  \node[c3,font=\small,align=center] at (axis cs:790,28) {$\Delta_rG^\circ>0$\\$K^\circ<1$};
  \node[font=\small,anchor=west] at (axis cs:310,80) {pente $= -\Delta_rS^\circ = +\SI{0,198}{kJ.K^{-1}.mol^{-1}}$};
\end{axis}
\end{tikzpicture}
\caption{$\Delta_rG^\circ(T)$ pour \ce{N2 + 3H2 = 2NH3} dans l'approximation d'Ellingham ($\Delta_rH^\circ = \SI{-91,8}{kJ.mol^{-1}}$, $\Delta_rS^\circ = \SI{-198,1}{J.K^{-1}.mol^{-1}}$).}
\label{fig:nh3drg}
\end{figure}

\begin{definition}[température d'inversion]
La \textbf{température d'inversion} $T_i$ est la température où $\Delta_rG^\circ(T_i) = 0$, soit $K^\circ(T_i) = 1$. Dans l'approximation d'Ellingham, $T_i = \Delta_rH^\circ/\Delta_rS^\circ$.
\end{definition}

\begin{exemple}[ammoniac, suite]
$\Delta_rG^\circ(\SI{298}{K}) = -91{,}8 - 298\times(-0{,}1981) = \SI{-32,7}{kJ.mol^{-1}}$ et $T_i = \dfrac{-91{,}8\times10^3}{-198{,}1} = \SI{463}{K}$ ($\SI{190}{\celsius}$). La synthèse est favorisée ($K^\circ > 1$) en dessous de $T_i$ : le terme enthalpique favorable domine à basse température, le terme entropique défavorable domine à haute température.
\end{exemple}

\begin{ordres}
\begin{itemize}
  \item $|\Delta_rH^\circ|$ : quelques dizaines à quelques centaines de $\mathrm{kJ\cdot mol^{-1}}$ (combustions : $\sim\SI{-800}{kJ.mol^{-1}}$ pour \ce{CH4}).
  \item $|\Delta_rS^\circ|$ : de l'ordre de $\SI{100}{}$ à $\SI{200}{J.K^{-1}.mol^{-1}}$ par mole de gaz créée ou consommée.
  \item À $\SI{298}{K}$, $T|\Delta_rS^\circ| \sim 30$ à $\SI{60}{kJ.mol^{-1}}$ : le terme entropique ne l'emporte sur le terme enthalpique qu'à haute température ou pour des réactions peu énergétiques.
\end{itemize}
\end{ordres}
